\section{Key Stakeholders \& Bloc Positions}

\lead{The tension between national security and free trade at the WTO involves several actors, from states to multilateral institutions whose interests and capacities shape how the issue is framed and negotiated.}

\subsection{Major security-driven powers}

\subsubsection{United States of America}

The United States of America is the central security-driven actor in today’s trading system, especially since the reelection of Donald J. Trump as 47th President of the country. It has been historically the principal architect of the liberal trade order we have today, despite being the major power currently straining it.\footnote{Komninakas, E. (2025, December 6). Architect to Arsonist: The United States’ Use of the National Security Exception in the World Trade Organization. Columbia Undergraduate Law Review. \url{https://www.culawreview.org/current-events-2/architect-to-arsonist-the-united-states-use-of-the-national-security-exception-in-the-world-trade-organization}} Since 2017, successive U.S. administrations have invoked national security to justify a widening range of trade restrictions, including Section 232 tariffs on steel and aluminium, broader tariff surcharges, and sweeping export controls on advanced semiconductors and computing hardware. \footnote{Reinsch, W. A., Caporal, J., \& CSIS. (2019). The WTO’s First Ruling on National Security: What Does It Mean for the United States? [Review of The WTO’s First Ruling on National Security: What Does It Mean for the United States?]. Csis.Org. \url{https://www.csis.org/analysis/wtos-first-ruling-national-security-what-does-it-mean-united-states}} \footnote{Branstetter, L. (2024, July). EXPORT CONTROLS AND US-CHINA TECHNOLOGY COMPETITION IN AN INTERDEPENDENT WORLD. Economic Studies at Brookings. \url{https://www.brookings.edu/wp-content/uploads/2024/07/20240701_Branstetter_Sanctions}. pdf} These policies rely on a domestic legal framework, Section 232 of the Trade Expansion Act, that allows the President to restrict imports deemed to “threaten to impair” national security, and they are frequently defended under GATT Article XXI when challenged in Geneva.\footnote{Komninakas, E. (2025, December 6). Architect to Arsonist: The United States’ Use of the National Security Exception in the World Trade Organization. Columbia Undergraduate Law Review.} \footnote{ARTICLE XXI Security Exceptions. (n.d.). In Analytical index of the GATT (pp. 599–610). \url{https://www.wto.org/english/res_e/booksp_e/gatt_ai_e/art21_e.pdf}} From the WTO’s perspective the U.S. position is particularly disruptive because it combines extensive use of the security exception along with a deliberate weakening of enforcement, as since 2017, Washington has blocked all appointments and reappointments to the Appellate Body, preventing the formation of the three judge quorum required for the body to hear appeals and thereby consequently weakening the binding nature of dispute settlement.\footnote{CSIS. (2025). The World Trade Organization: The Appellate Body Crisis | Economics Program and Scholl Chair in International Business | CSIS. Csis.org. \url{https://www.csis.org/programs/economics-program-and-scholl-chair-international-busine} ss/world-trade-organization} When WTO panels ruled in 2022 that the 2018 steel and aluminium tariffs did not meet the “war or other emergency in international relations” threshold, the United States rejected the decisions and argued that WTO bodies have no authority to review national-security measures at all, directly contradicting earlier jurisprudence such as Russia – Traffic in Transit.\footnote{Komninakas, E. (2025, December 6). Architect to Arsonist: The United States’ Use of the National Security Exception in the World Trade Organization. Columbia Undergraduate Law Review. \url{https://www.culawreview.org/current-events-2/architect-to-arsonist-the-united-states-use-of-the-national-security-exception-in-the-world-trade-organization}} \footnote{Sharon, J., \& Pace International Law Review. (2026, April 3). When “National Security” Becomes a Trade Weapon [Review of When “National Security” Becomes a Trade Weapon]. Pace International Law Review. \url{https://pilr.blogs.pace.edu/2026/04/03/when-national-security-becomes-a-trade-weapon/}} Ultimately, Washington has been using a combination of broad security claims, aggressive tariff policy along with institutional veto which makes it both a rule-taker and rule-breaker. It signals to other powerful states the possibility of insulating themselves from multilateral constraint if security is invoked. \footnote{Bacchus, J. (2022). The Black Hole of National Security Striking the Right Balance for the National Security Exception in International Trade Article XXI of the General Agreement on Tariffs and Trade (GATT) reserves for members of the [Review of The Black Hole of National Security Striking the Right Balance for the National Security Exception in International Trade Article XXI of the General Agreement on Tariffs and Trade (GATT) reserves for members of the]. \url{https://www.cato.org/sites/cato.org/files/2022-11/policy-analysis-936.pdf}}

\url{https://www.culawreview.org/current-events-2/architect-to-arsonist-the-united-states-use-of-the-national-security-exception-in-the-world-trade-organization}\href{https://www.culawreview.org/current-events-2/architect-to-arsonist-the-united-states-use-of-the-national-security-exception-in-the-world-trade-organization}{of-the-national-security-exception-in-the-world-trade-organization} People’s Republic of China

China is the other major pole in the emerging security-trade debate and its position is quite complex, being both a target of Western restrictions while becoming an increasingly proactive user of security-framed trade instruments.\footnote{FEI Xiuyan. (2022). National Security Considerations in China’s Trade Legislations: Offensive or Defensive? US-China Law Review, 19(4). \url{https://doi.org/10.17265/1548-6605/2022.04.004}} \footnote{Zhou, Weihuan and Jiang, Huiqin and Chen, Zhe, Trade vs. Security: Recent Developments of Global Trade Rules and China's Policy and Regulatory Responses from Defensive to Proactive (June 1, 2022). Forthcoming the World Trade Review, UNSW Law Research Paper No. 22-12, Available at SSRN: \url{https://ssrn.com/abstract=4131782}} On one hand, Chinese firms have been directly hit by U.S. and allied export controls on advanced chips, lithography equipment and many other technologies.\footnote{Sutter, K. (2025). U.S. Export Controls and China: Advanced Semiconductors [Review of U.S. Export Controls and China: Advanced Semiconductors]. Congress.Gov. \url{https://www.congress.gov/crs-product/R48642}} On the other hand, China itself has responded by broadening its own conception of economic security, introducing export controls on key minerals and technologies and developing a domestic legal framework that traits a wide range of economic activities as potential security concerns.\footnote{Andersen Institute. (2026, April 29). China’s Export Controls: Critical Minerals and Strategic Pressure Points. Andersen Institute. \url{https://anderseninstitute.org/chinas-export-control-architecture-and-its-use-of-critical-mi} nerals-as-strategic-pressure-points/} Recent analyses show that China’s strategy has shifted from largely defensive to more proactive: internationally, Beijing seeks to shape the interpretation of security exceptions and resist what it sees as abusive Western invocations; domestically, it has expanded national security to cover critical infrastructure, data, technology supply chains and industrial competitiveness. In WTO fora, China has repeatedly accused the U.S. of using national security as a cover for discriminatory Section 301 tariffs and industrial policy. In the meantime, Western members criticized China’s industrial subsidies and lack of transparency on state support, which creates a mutual securitisation of trade where each side is treating the other’s measures as illegitimate while defending its own. Therefore, this tension is key to any discussion on the future of Article XXI and the viability of the WTO as a neutral arbiter.

\subsection{Reform-oriented middle powers}

\subsubsection{The European Union}

The EU positions itself as the main champion of a fully functioning and legally predictable dispute settlement system even if it acknowledges the political difficulties of restoring the Appellate Body. Furthermore, EU policy documents underline that the WTO’s dispute settlement is a “central pillar” of the multilateral trading system and that its paralysis undermines the enforceability of rules and confidence in negotiations. In response, the EU has both pushed for Dispute Settlement Understanding (DSU) reform, through proposals on Appellate Body appointments, timelines, and institutional accountability, and developed instruments such as the Enforcement Regulation to deter and respond to partners that block effective dispute settlement.

One of the EU’s most important strategies is the Multi-Party Interim Appeal Arbitration Arrangement (MPIA), launched with other members under Article 25 of the DSU as a substitute for the Appellate Body. The MPIA preserves a two-tier system between participating members by providing for appeal arbitration by a panel of independent experts, an element that would reduce the incentive to appeal “into the void”.\footnote{World Trade Organization. (2020). WTO | Alternative Dispute Resolution Procedures [Review of WTO | Alternative Dispute Resolution Procedures]. Wto.Org. \url{https://www.wto.org/english/tratop_e/dispu_e/altds_e.htm}} Despite those efforts, the EU has repeatedly stressed that the MPIA is temporary and that the ultimate goal remains to restore a multilateral appellate mechanism, while inviting other members to join the arrangement to protect rules-based trade while negotiations are ongoing.\footnote{Machado, J. A. (2024). EU Statement at the General Council Meeting, 17 December 2024 [Review of EU Statement at the General Council Meeting, 17 December 2024]. EEAS. \url{https://www.eeas.europa.eu/delegations/world-trade-organization-wto/eu-statement-gen} eral-council-meeting-17-december-2024\_en} On the substantive side, the EU is not against national-security measures, but is concerned about the abuse of the security rhetoric in order to bypass safeguard disciplines. The EU has attacked the findings of the panels for what it thinks are an underestimation of the safeguard character of certain U.S. measures, insisting that international law be the only basis of WTO arbitration, not unilateral domestic definitions like “national security.”

\subsubsection{Japan}

Japan is a highly trade-dependent economy that has relied heavily on the WTO’s dispute settlement system and has become a key voice for its revival and reform (Ministry of Foreign Affairs of Japan, 2019). Japanese policymakers have said publicly that the WTO is “indispensable” for a trading nation, and have cited the loss of credibility and practical remedies that result from the Appellate Body’s paralysis. At the G20 and WTO, Japan has advocated for addressing concerns about the Appellate Body’s overreach, such as delayed rulings and perceived failure to fully resolve the measure at issue, while preserving the core features of a binding, independent appellate system.

After initial hesitation, the country finally decided to join the MPIA in March 2023, framing this choice as a way to maintain predictability and strengthen the rules-based multilateral trading system.\footnote{Australian National University. (2024). Japan’s joining MPIA an outside chance to boost momentum for WTO reform | Crawford School of Public Policy [Review of Japan’s joining MPIA an outside chance to boost momentum for WTO reform | Crawford School of Public Policy]. Crawford School of Public Policy. \url{https://crawford.anu.edu.au/australia-japan-research-centre/content-centre/article/news/ja} pans-joining-mpia-outside-chance} The cabinet understanding on MPIA participation states that the arrangement will substitute, on an interim basis, for the Appellate Body’s functions and that participation reflects Japan’s belief in the importance of an effective dispute-settlement framework.

At the same time, Japan is navigating domestic debates over economic security, especially regarding critical infrastructure such as high technologies and AI. To manage this, Tokyo is pursuing a policy of selective openness aiming to balance its need for foreign direct investments and technological innovation with a comprehensive screening framework designed to protect their domestic supply chains from foreign interference and espionage.\footnote{Edelman. (2026, June 2). Strengthening Japan’s Intelligence and Economic Security [Review of Strengthening Japan’s Intelligence and Economic Security]. Edelman. \url{https://www.edelman.com/insights/strengthening-japans-intelligence-and-economic-sec} urity} Ultimately, when it comes to the WTO, Japan will likely support proposals that rebuild enforceable rules with carefully designed reforms, rather than either dismantling legal review or leaving security claims entirely unconstrained.

\subsubsection{Canada}

Canada plays a role of great importance among reform-oriented middle powers most visibly through its leadership of the Ottawa Group on WTO Reform.\footnote{Global Affairs Canada. (2018, October 25). Joint Communiqué of the Ottawa Ministerial on WTO Reform. Canada.Ca; Government of Canada. \url{https://www.canada.ca/en/global-affairs/news/2018/10/joint-communique-of-the-ottawa-m} inisterial-on-wto-reform.html} This group of like-minded members composed of the EU, Japan, the United Kingdom, Australia, Mexico among others, aims to brainstorm pragmatic solutions and ideas to strengthen the three functions the WTO covers, from negotiations, to deliberation without forgetting dispute settlement. The latter is considered in joint ministerial statements from Ottawa as a central pillar of the organization, eventually warning that continued vacancies in the Appellate Body threaten the system as a whole. Canada has also been closely involved in early bilateral and plurilateral arrangements that paved the way for the MPIA, and is itself a participant in that interim appellate system. Canadian policy documents and parliamentary research highlight three main priorities: safeguarding and strengthening dispute settlement, improving transparency and notifications, and updating rules to address twenty-first century issues such as digital trade and industrial policy.

Ottawa’s stance is therefore one of multilateralism and reform as it accepts that some members have legitimate concerns about how dispute settlement has operated but insists that solutions must reinforce not weaken enforceability and transparency.

\subsection{Other emerging powers}

\subsubsection{India}

India is a large emerging market and power whose trade and foreign policy are directly articulated around strategic autonomy in the sense of being able to engage in free trade with all major powers while preserving independent decision making on security, energy and development.\footnote{Drishti IAS. (2025). Strategic Autonomy as India’s Global Compass [Review of Strategic Autonomy as India’s Global Compass]. Drishti IAS. \url{https://www.drishtiias.com/daily-updates/daily-news-editorials/strategic-autonomy-as-indi} a-s-global-compass} The country is effectively balancing a legacy of protectionism and a strong defense of the Global South with a rapid, proactive pivot toward deep global integration marked by free trade agreements that removed or deeply reduced tariffs on 96\% of trade with the European Union.\footnote{Dash, S. (2025, July 26). India tests strategic autonomy in fractured trade order. East Asia Forum. \url{https://eastasiaforum.org/2025/07/26/india-tests-strategic-autonomy-in-fractured-trade-or} der/} \footnote{Kuntikanamata, K. (2026, February 8). India’s trade deals test limits of strategic autonomy | Policy Circle. Policy Circle. \url{https://www.policycircle.org/opinion/india-eu-trade-deal/}} In the debate over national security versus free trade, India rejects both Western-led sweeping security restrictions and China's dominance, opting instead for a doctrine of Strategic Multialignment.\footnote{Drishti IAS. (2025). Strategic Autonomy as India’s Global Compass [Review of Strategic Autonomy as India’s Global Compass]. Drishti IAS. \url{https://www.drishtiias.com/daily-updates/daily-news-editorials/strategic-autonomy-as-indi} a-s-global-compass} India strongly opposes unilateral trade measures that are disguised as national security protections. When the United States enacted Section 232 tariffs on steel and aluminum under a national security banner, India was among the key nations that filed disputes at the WTO.\footnote{Reinsch, W. A., Caporal, J., \& CSIS. (2019). The WTO’s First Ruling on National Security: What Does It Mean for the United States? [Review of The WTO’s First Ruling on National Security: What Does It Mean for the United States?]. Csis.Org. \url{https://www.csis.org/analysis/wtos-first-ruling-national-security-what-does-it-mean-united-states}} New Delhi argues that Article XXI is not completely self-judging. India maintains that WTO panels must have the right to objectively review security claims to ensure developed nations do not use "national security" as a loophole for naked economic protectionism. Moreover, Modi’s government has criticized the EU’s Carbon Border Adjustment Mechanism (CBAM) and similar Western green-subsidies. New Delhi frames these policies as protectionist barriers that unfairly penalize developing economies under the guise of global climate security.\footnote{Mishra, A. (2026, June 2). Carbon at the border: India must turn its climate-trade argument into architecture. Down To Earth. \url{https://www.downtoearth.org.in/climate-change/carbon-at-the-border-india-must-turn-its-climate-trade-argument-into-architecture}}

While defending its policy space at the WTO, India has fundamentally shifted its economic strategy from defensive isolation to several free trade agreements. Breaking a decade-long lull in trade diplomacy, India concluded monumental trade pacts, including the historic India-EU FTA and a comprehensive trade framework with the United States. New Delhi has also sealed major trade agreements with the UK, Oman, New Zealand, and the European Free Trade Association (EFTA).\footnote{Solanki, V., \& IISS. (2019). India’s renewed focus on free-trade agreements [Review of India’s renewed focus on free-trade agreements]. IISS. \url{https://www.iiss.or}} Such agreements represent a deliberate effort by India to position itself as a credible alternative to China for global manufacturing.

Ultimately, India will fiercely advocate for the Policy Space for Development, arguing that developing nations must retain their right to protect their domestic agricultural and industrial activities while advocating for further integration of global supply chains.

\subsubsection{Russian Federation}

Following Russia’s invasion of Ukraine in 2022, the country has faced unparalleled, sweeping trade sanctions and the suspension of its Most-Favored-Nation (MFN) status by a coalition of roughly 40 nations including the G7 and the EU.\footnote{Kawashima, F. (2024). Trade Sanctions against Russia and their WTO Consistency: Focusing on Justification under National Security Exceptions. International Community Law Review, 26(1–2), 121–150. \url{https://doi.org/10.1163/18719732-12341497}} In response, Moscow has leveraged a dual strategy: utilizing WTO jurisprudence to defend its actions, while systematically decoupling from Western-led trade systems to build a parallel economic reality centered on "friendly countries." Nonetheless, long before the current wave of global sanctions, Russia fundamentally reshaped international trade law by forcing the WTO to rule on national security for the first time in its history.\footnote{Reinsch, W. A., Caporal, J., \& CSIS. (2019). The WTO’s First Ruling on National Security: What Does It Mean for the United States? [Review of The WTO’s First Ruling on National Security: What Does It Mean for the United States?]. Csis.Org. \url{https://www.csis.org/analysis/wtos-first-ruling-national-security-what-does-it-mean-united-states}} Following political turmoil in 2014, Russia severely restricted Ukrainian goods transiting through Russian territory to Central Asia, invoking the national security exception under GATT Article XXI(b)(iii) which allows trade violations during a "war or other emergency in international relations" (Taira, 2024). Russia argued that its national security decisions were entirely self-judging and that the WTO had no right to review them.\footnote{Voon, Tania, Russia - Measures Concerning Traffic in Transit (September 09, 2019). Pre-edited, pre-peer review draft in American Journal of International Law (Forthcoming), Available at SSRN: \url{https://ssrn.com/abstract=3475219} or \url{http://dx.doi.org/10.2139/ssrn.3475219}} While the WTO initially rejected this absolute immunity, it ultimately ruled in Moscow’s favor, agreeing that a legitimate “emergency in international relations” existed. This decision on Russia - Traffic in Transit, handed a powerful legal precedent to nations looking to weaponize trade restrictions under the cover of national security.\footnote{Voon, Tania, Russia - Measures Concerning Traffic in Transit (September 09, 2019). Pre-edited, pre-peer review draft in American Journal of International Law (Forthcoming), Available at SSRN: \url{https://ssrn.com/abstract=3475219} or \url{http://dx.doi.org/10.2139/ssrn.3475219}}

At the WTO, Russia forcefully maintains that Western sanctions, asset freezes and export controls are blatant violations of international trade law while Putin directly accused such sanctions to negatively impact the global economy, eroding trust in the dollar and euro and accelerating a shift toward a multipolar world (Associated Press, 2026). Moscow argues that Western nations are using "national security" as a political shield to wage economic warfare. Russian diplomats argue that suspending its MFN status destroys the core foundation of the multilateral trading system.

Facing these sanctions, Russia has stopped trying to repair its relations with traditional partners and has instead pushed for state-directed trade policy to build localized insulation. Moscow has legalized "parallel imports" (allowing the unauthorized import of branded Western goods through third-party nations) and heavily rely on shadow fleets and non-Western financial rails (like the Central Bank of Russia's SPFS network) to bypass G7 price caps and SWIFT bans. The Russian government has mandated that at least 80\% to 86\% of its total exports must be redirected to designated "friendly countries", primarily China, India, Belarus, and Kazakhstan. Through mega-projects like the International Cooperation and Export initiative, Russia is aggressively integrating its trade into transport corridors running through the Baltic, Black, and Pacific seas directly into Asia, Latin America, and Africa, explicitly cutting out transit hubs controlled by Western nations (Russian Government, 2026).

Russia will likely hold an aggressive, unyielding stance against Western economic dominance, heavily leaning on the legal precedent set by Russia - Traffic in Transit to claim their absolute sovereign right to protect their security interests. Russia would also be the vanguard of a group of non-Western powers that rejects the “weaponized dollar”.

\subsubsection{Rest of BRICS+}

Following its historic expansions, the expanded BRICS+ bloc (comprising Brazil, Russia, India, China, South Africa, Egypt, Ethiopia, Iran, and the United Arab Emirates, alongside a broad ring of newly minted partner states) has emerged as the premier institutional platform for the Global South.\footnote{Garcia Herrero, A. (2024, November 25). BRICS is becoming a more solid construction [Review of BRICS is becoming a more solid construction]. Bruegel | The Brussels-Based Economic Think Tank. \url{https://www.bruegel.org/first-glance/brics-becoming-more-solid-construction}} While individual members hold wildly divergent security agendas, the bloc is unified by a single core directive: insulating the developing world from Western economic warfare and the weaponization of the global financial system.\footnote{FTI Consulting. (2026, March 26). BRICS 2026: Strategic Autonomy and Economic Integration in a Fragmenting Global Order - FTI Strategic Communications [Review of BRICS 2026: Strategic Autonomy and Economic Integration in a Fragmenting Global Order - FTI Strategic Communications]. FTI Strategic Communications. \url{https://fticommunications.com/brics-2026-strategic-autonomy-and-economic-integration-in-a-fragmenting-global-order/}} As codified in the watershed Kazan Declaration, BRICS+ officially condemns what it labels "unilateral coercive measures", specifically secondary sanctions, asset freezes, and Western export controls. The bloc formally maintains that these measures are fundamentally illegal under international law, damage the UN Charter, and actively undermine the rules-based multilateral trading system (Huashen, 2024). Led by India’s execution of Digital Public Infrastructure (DPI) and China’s cross-border payment integration, BRICS+ is actively expanding local currency settlement systems for intra-bloc trade. By trading in currencies like the Yuan, Rupee, and Dirham, members look to bypass the structural vulnerabilities of G7-dominated banking corridors (Zhou, 2024).

BRICS+ members would have a careful balance, as while presenting a united front against the expansion of Western “national security” definitions, they must also accurately mirror their specific domestic priorities.

\subsubsection{Global South and Least Developed Countries}

While major powers engage in geoeconomic competition, the broader Global South most notably represented by the G-90 coalition views the securitization of trade with existential alarm. Indeed, they lack the fiscal capacity needed to compete with Western-style green industrial subsidies, as well as the geopolitical leverage to fight weaponized tariff surcharges. Consequently, these nations are the primary victims of a fragmented multilateral trading order.\footnote{UNCTAD. (2025). Tariff disruptions [Review of Tariff disruptions]. \url{https://unctad.org/system/files/official-document/osgttinf2025d6_en.pdf}} In institutional frameworks like the African Union (AU) preparatory summits for WTO Ministerial Conferences, these countries have strongly condemned the rising use of national security exceptions as a convenient shield for raw protectionism. From a structural perspective, the ongoing paralysis of the WTO Appellate Body poses an immediate threat to low-income economies. Without a functioning, two-tier dispute settlement mechanism, smaller trading nations lose their only mechanism to defend their trade rights on an equal footing against economically dominant powers.\footnote{UNCTAD. (2025). Tariff disruptions [Review of Tariff disruptions]. \url{https://unctad.org/system/files/official-document/osgttinf2025d6_en.pdf}}

G-90 countries would position themselves as protectors of the strict, original text of the Marrakesh Agreement, emphasizing that the rule of law is the weapon of the weak while the law of the jungle only benefits the strong.
